% Modelo de artigo do I3exiga IFSP.
% Compile com: pdflatex modelo-artigo.tex
% Preencha os campos em "Dados do artigo" e escreva o texto nas seções abaixo.
% Artigos de SOLUÇÃO devem informar o link do problema; artigos de CONHECIMENTO GERAL
% não podem resolver problemas da plataforma. Códigos devem estar em C, C++, Java ou Python.
\documentclass[12pt,a4paper]{article}

\usepackage[utf8]{inputenc}
\usepackage[T1]{fontenc}
\usepackage[brazil]{babel}
\usepackage[margin=2.5cm]{geometry}
\usepackage{hyperref}
\usepackage{listings}
\usepackage{xcolor}

\lstset{
  basicstyle=\ttfamily\small,
  breaklines=true,
  frame=single,
  numbers=left,
  numberstyle=\tiny,
  showstringspaces=false
}

% ===================== Dados do artigo =====================
\newcommand{\tipoartigo}{Conhecimento geral} % ou: Solução de problema
\newcommand{\titulo}{Título do artigo}
\newcommand{\autores}{Nome do(a) autor(a) 1 \and Nome do(a) autor(a) 2}
\newcommand{\linkproblema}{} % apenas para artigos de solução, ex.: https://exemplo.ifsp.edu.br/problems/view?id=1
% ===========================================================

\title{\titulo}
\author{\autores}
\date{\today}

\begin{document}
\maketitle

\begin{center}
  \textbf{Tipo:} \tipoartigo \\
  \ifx\linkproblema\empty\else \textbf{Problema:} \url{\linkproblema} \fi
\end{center}

\begin{abstract}
Resumo do artigo em poucas linhas.
\end{abstract}

\section{Introdução}
Apresente o tema e o objetivo do artigo.

\section{Desenvolvimento}
Explique o conteúdo. Em artigos de solução, descreva a lógica de cada solução de forma explícita
antes de apresentar o código.

\begin{lstlisting}[language=C, caption={Exemplo de código em C}]
#include <stdio.h>

int main(void) {
    int a, b;
    scanf("%d %d", &a, &b);
    printf("%d\n", a + b);
    return 0;
}
\end{lstlisting}

\section{Conclusão}
Encerre com os principais pontos.

\begin{thebibliography}{9}
\bibitem{ref1} Sobrenome, N. \textit{Título da obra}. Editora, ano.
\end{thebibliography}

\end{document}
